\section{模型假设}

为简化问题并保证模型的合理性与可复现性，本文作出如下假设：

\begin{enumerate}[leftmargin=2.2em,itemsep=2pt]
  \item 假设数据集中图像的标注准确可靠，标注框能够完整覆盖缺陷区域，且类别标签与缺陷实际外观一致；
  \item 假设训练集与测试集图像均来自同一拍摄场景与采集设备，图像分辨率统一为 $640\times640$，满足同分布假设；
  \item 假设检测网络输出的置信度分数能够单调反映目标存在的真实概率，即置信度越高，该位置存在缺陷的可能性越大；
  \item 假设有缺陷图像中至少存在一个高置信度检测结果，而无缺陷图像的最大检测置信度显著偏低，二者在最大置信度统计量上可分；
  \item 假设各缺陷类别之间互不重叠、独立标注，一张图像中允许同时存在多个类别的缺陷；
  \item 假设模型推理阶段的计算资源充足，允许使用 GPU 加速以满足实时检测需求。
\end{enumerate}
